% Mathématiques Terminale spécialité — V3 LaTeX
% Représentations paramétriques et équations cartésiennes

\section{Objectifs}
\begin{itemize}
\item Déterminer et reconnaître une représentation paramétrique de droite.
\item Déterminer et reconnaître une équation cartésienne de plan.
\item Étudier des intersections dans l’espace.
\item Traduire une configuration géométrique par un système d’équations.
\end{itemize}

\section{Repères et enjeux du chapitre}
Les représentations paramétriques et les équations cartésiennes sont deux langages complémentaires. La droite est naturellement décrite par un point, une direction et un paramètre réel ; le plan est naturellement décrit par un point et un vecteur normal, ce qui conduit à une équation cartésienne. Une configuration géométrique devient alors un système algébrique dont les solutions traduisent les intersections.

\section{1. Représentation paramétrique d’une droite}
La droite passant par $A(x_A,y_A,z_A)$ et de vecteur directeur $\vec u=(a,b,c)$ admet la représentation
\[
\left\{\begin{array}{l}
x=x_A+at,\\
y=y_A+bt,\\
z=z_A+ct,
\end{array}\right.
\qquad t\in\mathbb R.
\]
Un seul paramètre $t$ commande simultanément les trois coordonnées. Plusieurs représentations paramétriques peuvent décrire la même droite, car on peut choisir un autre point de la droite ou multiplier le vecteur directeur par un réel non nul.

\section{2. Appartenance à une droite}
Pour vérifier qu'un point $M(x_M,y_M,z_M)$ appartient à une droite paramétrée, il faut trouver une même valeur de $t$ satisfaisant les trois équations. Une valeur compatible avec deux coordonnées mais pas avec la troisième prouve que le point n'appartient pas à la droite. Cette vérification simultanée évite une erreur fréquente liée à la projection de la configuration sur seulement deux axes.

\coursefigure{/images/mathematiques/terminale-specialite-mathematiques/representations-parametriques-equations-cartesiennes-1.svg}{Schéma structurant pour Représentations paramétriques et équations cartésiennes.}{La figure met en évidence les objets, relations ou variations fondamentales mobilisés dans la première moitié du chapitre.}

\section{3. Équation cartésienne d’un plan}
Si $\vec n=(a,b,c)$ est un vecteur normal non nul à un plan $\mathcal P$, alors une équation cartésienne de ce plan est
\[
ax+by+cz+d=0.
\]
Si le plan passe par $A(x_A,y_A,z_A)$, on peut écrire directement
\[
a(x-x_A)+b(y-y_A)+c(z-z_A)=0.
\]
Les coefficients $a$, $b$, $c$ sont donc les coordonnées d'un vecteur normal au plan et non d'un vecteur directeur du plan.

\section{4. Vérifier l’appartenance à un plan}
Un point appartient au plan si ses coordonnées satisfont son équation. La vérification consiste donc à substituer les trois coordonnées. Si le résultat est nul, le point est dans le plan. Cette opération est simple, mais elle devient un outil stratégique dans de nombreux exercices : déterminer la constante $d$, vérifier une intersection ou construire un plan passant par des points donnés.

\section{5. Intersection d’une droite et d’un plan}
On remplace les coordonnées paramétrées de la droite dans l'équation du plan. On obtient une équation en $t$. Une solution unique fournit un point d'intersection. Une égalité impossible correspond à une droite parallèle au plan et extérieure. Une identité correspond à une droite incluse dans le plan.
\[
ax(t)+by(t)+cz(t)+d=0.
\]
La nature géométrique de l'intersection se lit donc dans la nature de l'équation obtenue.

\coursefigure{/images/mathematiques/terminale-specialite-mathematiques/representations-parametriques-equations-cartesiennes-2.svg}{Deuxième représentation pédagogique pour Représentations paramétriques et équations cartésiennes.}{La figure sert de support de lecture, de comparaison ou de contrôle pour les méthodes centrales du chapitre.}

\section{6. Intersection de deux droites}
Pour deux droites paramétrées par $t$ et $s$, on égalise les trois coordonnées et on résout le système. Une solution $(t,s)$ commune fournit le point d'intersection. Si les vecteurs directeurs sont colinéaires, il faut distinguer parallélisme et confusion. S'ils ne sont pas colinéaires mais que le système n'a pas de solution, les droites sont non coplanaires.

\section{7. Construire un projeté orthogonal}
Pour projeter un point $M$ sur un plan $\mathcal P$, on construit la droite $d$ passant par $M$ et dirigée par un vecteur normal $\vec n$ à $\mathcal P$. On écrit une représentation paramétrique de $d$, puis on résout son intersection avec $\mathcal P$. Le point obtenu est le projeté orthogonal $H$. La démarche combine donc vecteur normal, paramétrage et équation cartésienne dans une seule stratégie.


\section{Méthode générale de résolution}
\begin{methode}
\begin{enumerate}
\item Identifier précisément l'objet mathématique demandé et les hypothèses disponibles.
\item Traduire l'énoncé dans une écriture adaptée : égalité, système, représentation paramétrique, inégalité, suite ou fonction selon le contexte.
\item Choisir le théorème ou la propriété en vérifiant explicitement ses conditions d'application.
\item Effectuer les calculs en conservant les formes exactes aussi longtemps que possible.
\item Contrôler la cohérence du résultat par une seconde méthode, un ordre de grandeur, un signe, une substitution ou une lecture graphique.
\item Conclure dans le langage du problème, en distinguant clairement ce qui est démontré de ce qui est conjecturé ou approché numériquement.
\end{enumerate}
\end{methode}

\section{Rédaction attendue en Terminale}
En spécialité, une réponse correcte ne se réduit pas à une ligne de calcul. Lorsqu'un résultat dépend d'un théorème, les hypothèses doivent apparaître dans la copie. Lorsqu'une valeur est approchée, la précision doit être annoncée. Lorsqu'une équation est résolue, il faut revenir à la question posée et vérifier que la solution appartient au domaine considéré. Cette discipline de rédaction évite les raisonnements implicites et prépare aux attendus de l'épreuve écrite.

\begin{attention}
Une calculatrice ou un logiciel peut suggérer une réponse, produire un tableau ou une représentation, mais il ne remplace pas la justification. Un résultat numérique devient exploitable seulement après avoir été relié à une propriété mathématique et interprété dans le contexte.
\end{attention}

\section{Contrôler une solution}
Le contrôle dépend du chapitre : recompter par le complément en combinatoire, vérifier l'appartenance d'un point à une droite ou à un plan, substituer une solution dans une équation, comparer deux expressions de limite, vérifier un signe de dérivée, ou encore contrôler qu'un encadrement de dichotomie conserve bien le changement de signe. Dans tous les cas, l'objectif est d'obtenir une preuve locale que le résultat final n'est pas seulement plausible mais compatible avec les données et les hypothèses.

\section{Problème de synthèse et stratégie}

Dans ce chapitre, la difficulté d'un problème complet consiste à articuler plusieurs résultats plutôt qu'à appliquer une formule isolée. Pour représentations paramétriques et équations cartésiennes, on commence par identifier les objets et les hypothèses, puis on construit une chaîne de décisions vérifiables. Les résultats intermédiaires doivent être réutilisés explicitement : une relation obtenue peut justifier une appartenance, une monotonie, un comptage, une limite ou un encadrement. Cette organisation est particulièrement importante lorsque l'énoncé introduit une notation nouvelle ou demande une interprétation finale.


\section{Réinvestissement type baccalauréat}
Un exercice de baccalauréat combine souvent plusieurs compétences : lecture d'une situation, choix d'une stratégie, calcul exact, justification et interprétation. Il faut donc savoir passer d'une question technique courte à une question de synthèse. Une bonne stratégie consiste à repérer les résultats intermédiaires qui seront réutilisés, à annoncer les théorèmes avant leur emploi et à conserver une trace claire des dépendances entre les questions. Lorsqu'une question paraît isolée, elle prépare souvent une propriété utile pour la suite de l'exercice.

\section{Erreurs fréquentes}
\begin{itemize}
\item Utiliser plusieurs paramètres indépendants pour une même droite.
\item Interpréter $(a,b,c)$ comme un vecteur directeur du plan $ax+by+cz+d=0$.
\item Conclure à une intersection de deux droites après seulement deux égalités de coordonnées.
\item Oublier qu’une droite parallèle à un plan peut aussi être incluse dans ce plan.
\item Confondre une équation cartésienne de plan avec une équation de droite dans l’espace.
\end{itemize}

\section{À retenir}
\begin{aretenir}
\begin{itemize}
\item Une droite se paramètre avec un seul réel.
\item Un plan d’équation $ax+by+cz+d=0$ a pour vecteur normal $(a,b,c)$.
\item L’appartenance se vérifie par substitution.
\item L’intersection droite-plan se ramène à une équation sur le paramètre.
\item Un système sans solution peut traduire des droites non coplanaires.
\end{itemize}
\end{aretenir}
